\bookchapter{บทที่ 30:การแคสต์หน้าจอ การนำเสนอ และการเชื่อมต่อระบบรถยนต์}

\chapterintro{หน้าจอมือถือถึงแม้จะใหญ่แค่ไหนก็แค่ 7 นิ้ว แต่ทีวีของคุณมี 55 นิ้ว โปรเจคเตอร์ของคุณมี 100 นิ้ว การ “โยน” ภาพจากมือถือไปยังหน้าจอขนาดใหญ่ ประสบการณ์ที่ได้จะอยู่คนละมิติเลย บทนี้จะอธิบายให้ชัดเจนเกี่ยวกับการทำงานของการแคสต์หน้าจอ การนำเสนอ และการเชื่อมต่อรถยนต์ ซึ่งเป็น 3 สถานการณ์ที่พบบ่อย เพื่อช่วยให้คุณขยายความสามารถของมือถือไปสู่หน้าจอที่ใหญ่ขึ้น}

\booksection{3 ค่ายการแคสต์หน้าจอไร้สาย — อย่าสับสน}

การแคสต์หน้าจอไร้สายดูเหมือนจะเป็น “ภาพจากมือถือไปที่ทีวี” แต่เทคโนโลยีเบื้องหลังมี 3 แบบ ซึ่งไม่สามารถใช้ร่วมกันได้:

\begin{center}
\begin{tabularx}{\linewidth}{|X|X|X|}
\hline
\textbf{เทคโนโลยี} & \textbf{ใครใช้} & \textbf{จุดเด่น} \tabularnewline
\hline
AirPlay & iPhone/iPad/Mac & ระบบนิเวศ Apple ประสบการณ์ลื่นไหลที่สุด ต้องใช้ Apple TV หรือทีวีที่รองรับ AirPlay \tabularnewline
Miracast & มือถือ Android & เชื่อมต่อโดยตรงผ่าน Wi-Fi Direct ไม่ต้องใช้เราเตอร์ มือถือ Android จะใช้วิธีนี้เป็นค่าเริ่มต้นสำหรับ "การแคสต์หน้าจอไร้สาย" \tabularnewline
DLNA & ทุกแพลตฟอร์ม & ส่งเฉพาะไฟล์สื่อ (วิดีโอ/เพลง/รูปภาพ) ไม่มิเรอร์หน้าจอ เหมาะสำหรับส่งวิดีโอไปเล่นบนทีวี \tabularnewline
\hline
\end{tabularx}
\end{center}

\booksubsection{วิธีการใช้งานของแต่ละแพลตฟอร์ม:}

\begin{itemize}
  \item \textbf{iPhone}:ศูนย์ควบคุม > การมิเรอร์หน้าจอ > เลือกทีวี/Apple TV ของคุณ หากทีวีไม่รองรับ AirPlay คุณสามารถซื้อกล่อง Apple TV หรือสติ๊กทีวีที่รองรับ AirPlay
  \item \textbf{Android}:เลื่อนแถบการแจ้งเตือนลง > ค้นหาปุ่ม "แคสต์หน้าจอไร้สาย"/"แคสต์หน้าจออัจฉริยะ"/"โต้ตอบหลายหน้าจอ" > เลือกทีวี ต้องเชื่อมต่อมือถือและทีวีเข้ากับ Wi-Fi เดียวกัน
  \item \textbf{โซลูชันสารพัดประโยชน์}:หากทีวีเก่าเกินไปและไม่รองรับการแคสต์หน้าจอแบบไร้สายใดๆ ให้ซื้อ Chromecast หรือไมโครกล่อง Xiaomi/กล่อง Huawei เสียบเข้ากับพอร์ต HDMI ของทีวี ก็จะสามารถรับสัญญาณแคสต์หน้าจอได้แล้ว
\end{itemize}

\booksection{การแคสต์หน้าจอแบบมีสาย — มีเสถียรภาพและไม่มีความละเหลี่ยมเวลา}

การแคสต์หน้าจอแบบไร้สายสะดวกแต่มีความล่าช้า ประสบการณ์การใช้งานเวลาเล่นเกมและนำเสนอจึงไม่ค่อยดีนัก การแคสต์หน้าจอแบบมีสายสามารถแก้ปัญหานี้ได้:

\begin{itemize}
  \item \textbf{iPhone}:ต้องใช้สายแปลง Lightning/USB-C เป็น HDMI (ราคาประมาณ 100-200 หยวน) เสียบแล้วหน้าจอมือถือจะส่งออกไปยังพอร์ต HDMI ของทีวี/โปรเจคเตอร์โดยตรง
  \item \textbf{Android}:มือถือบางรุ่นรองรับการเชื่อมต่อ USB-C กับ HDMI โดยตรง (ต้องมือถือรองรับ DP Alt Mode ซึ่ง flagship ส่วนใหญ่รองรับ) แค่ซื้อสาย USB-C to HDMI มาก็พอ
  \item \textbf{สถานการณ์ในห้องประชุม}:หากคุณต้องใช้มือถือนำเสนอบ่อยๆ แนะนำให้พกสายแปลงเป็น HDMI ติดกระเป๋าไว้เสมอ ทีวี/โปรเจคเตอร์ในห้องประชุมเกือบทุกที่จะมีพอร์ต HDMI เสียบก็ใช้ได้เลย น่าเชื่อถือกว่าการนั่งปวดหัวกับการแคสต์แบบไร้สายมาก
\end{itemize}

\booksection{เปลี่ยนมือถือเป็นคอมพิวเตอร์ — โหมดเดสก์ท็อป}

มือถือ Android บางรุ่นรองรับ “โหมดเดสก์ท็อป” — เมื่อต่อกับจอภาพแล้ว มือถือจะแสดงผลหน้าจอคล้ายกับคอมพิวเตอร์เดสก์ท็อป สามารถใช้งานร่วมกับเมาส์และคีย์บอร์ดได้:

\begin{itemize}
  \item \textbf{Samsung DeX}:ฟังก์ชันเด่นของ Samsung flagship หลังจากเชื่อมต่อกับจอภาพ หน้าจอเดสก์ท็อปจะปรากฏขึ้น สามารถใช้งานหลายหน้าต่างพร้อมกันได้ ประสบการณ์การใช้งานใกล้เคียงกับคอมพิวเตอร์จริงๆ มากที่สุด ร่วมกับคีย์บอร์ดและเมาส์ Bluetooth ตอนเดินทางก็ไม่ต้องพกโน้ตบุ๊กแล้ว
  \item \textbf{โหมดเดสก์ท็อปของ Huawei}:ซีรีส์ Huawei Mate/Pura หลังจากเชื่อมต่อกับจอภาพก็มีฟังก์ชันที่คล้ายกัน รองรับการใช้งานหลายหน้าต่าง การจัดการไฟล์ และ App สำนักงาน
  \item \textbf{มอเตอโรล่า Ready For}:Moto flagship ก็รองรับเช่นกัน มีฟังก์ชันคล้ายกับ DeX
  \item \textbf{คุณสมบัติ}:งานสำนักงานเบาๆ (ตอบอีเมล แก้ไขเอกสาร) นำเสนอ PPT ดูวิดีโอบนจอใหญ่ ไม่เหมาะกับการใช้งาน Excel หนักๆ หรืองานออกแบบระดับมืออาชีพ
\end{itemize}

\booksection{การเชื่อมต่อรถยนต์ — วิธีที่ถูกต้องในการใช้มือถือขณะขับรถ}

การใช้มือถือขณะขับรถเป็นสิ่งผิดกฎหมาย แต่การนำทาง ฟังเพลง และรับสายกลับเป็นสิ่งจำเป็น ระบบการเชื่อมต่อรถยนต์ช่วยให้คุณใช้ฟังก์ชันหลักของมือถือบนหน้าจอระบบรถยนต์ได้ ส่วนมือถือนั้นสามารถวางไว้ด้านข้างโดยไม่ต้องแตะ:

\begin{itemize}
  \item \textbf{Apple CarPlay}:ตัวเลือกแรกสำหรับผู้ใช้ iPhone เชื่อมต่อกับระบบรถยนต์ผ่าน USB หรือแบบไร้สาย หน้าจอระบบรถยนต์จะแสดงอินเทอร์เฟซแบบย่อของ iPhone — นำทาง (Amap/Baidu), เพลง, โทรศัพท์, ข้อความ, พอดแคสต์ ใช้ Siri ควบคุมด้วยเสียงได้ตลอดโดยไม่ต้องใช้มือ
  \item \textbf{Android Auto}:โซลูชันสำหรับผู้ใช้ Android มีฟังก์ชันคล้ายกับ CarPlay รองรับ Google Maps, Spotify ฯลฯ แต่ในประเทศไทยเนื่องจากข้อจำกัดของบริการ Google ฟังก์ชันบางส่วนอาจใช้งานไม่ได้
  \item \textbf{Huawei HiCar}:โซลูชันที่ Huawei พัฒนาเอง ได้รับการรองรับจากรถยนต์รุ่นต่างๆ ในประเทศไทยในอัตราที่สูง นอกจากนำทางและเพลงแล้ว ยังมีฟังก์ชันเด่นอย่างการโทรวิดีโอผ่าน Huawei และการผสานรวมภาพจากกล้องถอยรถอีกด้วย
  \item \textbf{Baidu CarLife}:โซลูชันข้ามแพลตฟอร์ม ทั้ง iOS และ Android ใช้ได้ ฟังก์ชันน้อยกว่าแต่ความเข้ากันได้กว้างขวาง รถรุ่นเก่าหลายคันรองรับเฉพาะ CarLife เท่านั้น
\end{itemize}

\begin{tcbwarn}{คำเตือนด้านความปลอดภัย}
ไม่ว่าจะเลือกโซลูชันการเชื่อมต่อรถยนต์แบบใด ขณะขับรถควรพยายามใช้การควบคุมด้วยเสียงให้มากที่สุด ตั้งจุดหมายนำทางล่วงหน้า เลือกเพลย์ลิสต์ล่วงหน้า และห้ามไปจิ้มหน้าจอระบบรถยนต์ขณะขับขี่

\end{tcbwarn}

\booksection{การแคสต์หน้าจอในห้องประชุมจริง — ไม่ต้องอลหม่านอีกต่อไป}

การประชุมเป็นสถานการณ์ที่ต้องใช้การแคสต์หน้าจอบ่อยที่สุด และก็เป็นสถานการณ์ที่มักมีปัญหามากที่สุดเช่นกัน นี่คือประสบการณ์จริงไม่กี่ข้อ:

\begin{itemize}
  \item \textbf{ไปถึงก่อนเวลาเพื่อทดสอบ}:อย่ารอจนถึงเวลาประชุมจึงค่อยต่อแคสต์หน้าจอ ไปก่อน 10 นาทีเพื่อทดสอบวิธีการเชื่อมต่อ อัตราความล้มเหลวของการแคสต์ผ่าน Wi-Fi นั้นสูงกว่าที่คุณคิด
  \item \textbf{เตรียมวิธีแบบมีสายไว้เป็น Plan B}:พกสายแปลงเป็น HDMI ไว้ในกระเป๋า เมื่อแคสต์แบบไร้สายไม่ได้ สายจะเป็นท่องช่วยชีวิตคุณ
  \item \textbf{ปิดการแสดงตัวอย่างการแจ้งเตือน}:เวลาแคสต์หน้าจอ ข้อความ WeChat ของคุณจะแสดงบนจอใหญ่ในห้องประชุม ก่อนแคสต์ให้ไปที่การตั้งค่าเพื่อปิด "การแสดงตัวอย่างเนื้อหาการแจ้งเตือนบนหน้าจอล็อก" หรือเปิด "โหมดการนำเสนอ" (มือถือบางรุ่นรองรับ จะซ่อนการแจ้งเตือนอัตโนมัติเมื่อแคสต์)
  \item \textbf{ใช้ "เอกสารนำเสนอ" แทน "การใช้งานมือถือ"}:หากต้องการนำเสนอ PPT ให้บันทึก PPT ลงในมือถือแล้วเปิดด้วย WPS เพื่อเล่นแบบเต็มจอ แทนที่จะแคสต์หน้าจอมือถือทั้งหมด วิธีนี้จะดูเป็นมืออาชีพและลื่นไหลกว่า
\end{itemize}

\begin{tcbpractice}{ลงมือทำ: ทำการแคสต์หน้าจอจากมือถือไปยังทีวีสำเร็จ}
\begin{enumerate}
  \item : ยืนยันว่ามือถือกับทีวีต่อ Wi-Fi เดียวกัน
  \item เปิดรูปภาพหรือวิดีโอในมือถือ
  \item ลองแคสต์แบบไร้สาย (iPhone ใช้การกระจกภาพหน้าจอ, Android ใช้ปุ่มแคสต์หน้าจอแบบไร้สาย)
  \item หากไม่สำเร็จ ให้ตรวจสอบว่าทีวีรองรับโปรโตคอลที่ตรงกันหรือไม่ หรือลองใช้ DLNA push (กดปุ่ม "แคสต์" ในแอปอัลบั้ม/วิดีโอ)
  \item หลังแคสต์สำเร็จแล้ว ลองทำงานอย่างอื่นบนมือถือ แล้วดูว่าหน้าจอใหญ่จะแสดงผลพร้อมกันหรือไม่
\end{enumerate}

\end{tcbpractice}

\begin{tcbsummary}{สรุปบทนี้}
\begin{itemize}
  \item สามค่ายของการแคสต์ไร้สาย: AirPlay (Apple), Miracast (Android), DLNA (ส่งสื่อข้ามแพลตฟอร์ม)
  \item การแคสต์แบบมีสายมีเสถียรภาพและไม่มีความล่าช้า ต้องมีสายแปลง HDMI สำหรับสถานการณ์ในห้องประชุม
  \item Samsung DeX/โหมดเดสก์ท็อปของ Huawei เปลี่ยนมือถือให้กลายเป็นคอมพิวเตอร์สำนักงานเบาๆ
  \item การเชื่อมต่อรถยนต์: โซลูชันหลัก 4 แบบ ได้แก่ CarPlay/Android Auto/HiCar/CarLife
  \item กฎเหล็กของการแคสต์ในห้องประชุม: ทดสอบล่วงหน้า, เตรียม Plan B, ปิดการแสดงตัวอย่างการแจ้งเตือน
\end{itemize}

\end{tcbsummary}
